\section{ابزارهای لازم}
برای این مرحله به ابزارهای زیر نیاز است:

\begin{itemize}
  \item
  \lr{VS Code}
  و افزونهٔ
  \lr{PlatformIO IDE}
  ؛
  \item کارگزار و ابزارهای خط فرمان
  \lr{Mosquitto}
  ؛
  \item
  \lr{TeXstudio}
  و
  \lr{XeLaTeX}
  برای ویرایش همین سند.
\end{itemize}

\subsection{نصب
\lr{Mosquitto}
}
در
\lr{Windows}
نصب‌کنندهٔ رسمی
\lr{Mosquitto}
را اجرا و مسیر نصب را به
\lr{PATH}
اضافه کنید. در
\lr{Linux}
یا
\lr{macOS}
از فرمان متناسب استفاده کنید:

\begin{LTR}
\begin{lstlisting}[language=bash]
# Ubuntu / Debian
sudo apt update
sudo apt install mosquitto mosquitto-clients

# macOS with Homebrew
brew install mosquitto
\end{lstlisting}
\end{LTR}

نصب را در یک پایانهٔ تازه بررسی کنید:

\begin{LTR}
\begin{lstlisting}
mosquitto -h
mosquitto_pub -h
mosquitto_sub -h
\end{lstlisting}
\end{LTR}

\section{آزمایش مستقل روی رایانه}
این آزمایش باید پیش از روشن‌کردن برد موفق شود.

\begin{stepbox}{گام ۱: اجرای کارگزار}
در پایانهٔ اول اجرا کنید:

\begin{LTR}
\begin{lstlisting}
mosquitto -v
\end{lstlisting}
\end{LTR}
\end{stepbox}

\begin{stepbox}{گام ۲: انتظار برای پیام}
در پایانهٔ دوم اجرا کنید:

\begin{LTR}
\begin{lstlisting}
mosquitto_sub -h localhost -t "iot/test" -v
\end{lstlisting}
\end{LTR}
\end{stepbox}

\begin{stepbox}{گام ۳: ارسال پیام}
در پایانهٔ سوم اجرا کنید:

\begin{LTR}
\begin{lstlisting}
mosquitto_pub -h localhost -t "iot/test" -m "broker-ok"
\end{lstlisting}
\end{LTR}
\end{stepbox}

خروجی مورد انتظار در پایانهٔ دوم:

\begin{LTR}
\begin{lstlisting}
iot/test broker-ok
\end{lstlisting}
\end{LTR}

اگر این خروجی دیده نشد، هنوز سراغ برد نروید. ابتدا اجرای سرویس، شمارهٔ درگاه و پیام خطای پایانهٔ اول را بررسی کنید.

\section{اجازهٔ اتصال از شبکهٔ محلی}
اتصال
\lr{localhost}
فقط روی همان رایانه کار می‌کند. برای آزمایش محلی برد، فایلی به نام
\lr{mosquitto-local.conf}
در یک پوشهٔ موقت بسازید:

\begin{LTR}
\begin{lstlisting}
listener 1883
allow_anonymous true
\end{lstlisting}
\end{LTR}

کارگزار آزمایشی را با این فایل اجرا کنید:

\begin{LTR}
\begin{lstlisting}
mosquitto -c mosquitto-local.conf -v
\end{lstlisting}
\end{LTR}

\begin{notebox}
این تنظیم فقط برای توسعه در شبکهٔ خصوصی و مطمئن است. درگاه
\lr{1883}
را روی روتر به اینترنت منتقل نکنید. در نسخهٔ نهایی، نام کاربری، رمز و ارتباط رمزگذاری‌شده اضافه خواهد شد.
\end{notebox}

\section{پیداکردن نشانی رایانه}
برد به نشانی
\lr{IPv4}
رایانه نیاز دارد. یکی از فرمان‌های زیر را اجرا کنید:

\begin{LTR}
\begin{lstlisting}
# Windows
ipconfig

# Linux
ip addr

# macOS
ifconfig
\end{lstlisting}
\end{LTR}

برد و رایانه باید به یک شبکه متصل باشند. در دیوارهٔ آتش، فقط دسترسی ورودی
\lr{1883/TCP}
از شبکهٔ خصوصی را مجاز کنید.
